\documentclass[aspectratio=169,10pt]{beamer}
\usetheme{Madrid}
\usecolortheme{default}

\usepackage[utf8]{inputenc}
\usepackage[T1]{fontenc}
\usepackage[brazilian]{babel}
\usepackage{booktabs}
\usepackage{amsmath}
\usepackage{xspace}

% Números gerados por paper/gera_tabelas.py a partir de sim/results/.
% Gerado por paper/gera_tabelas.py. Não editar à mão.
\newcommand{\NumMatrizes}{8\xspace}
\newcommand{\NumExecSpmv}{96\xspace}
\newcommand{\NumBitExatas}{96\xspace}
\newcommand{\MinLinhas}{30\xspace}
\newcommand{\MaxLinhas}{1033\xspace}
\newcommand{\MaxNnz}{6027\xspace}
\newcommand{\MedFPtrinta}{$5{,}91 \times 10^{-8}$\xspace}
\newcommand{\MedTakumTrinta}{$1{,}15 \times 10^{-8}$\xspace}
\newcommand{\MedFPtrintaAcSes}{$3{,}61 \times 10^{-8}$\xspace}
\newcommand{\MedTakumTrintaAcTrinta}{$2{,}94 \times 10^{-8}$\xspace}
\newcommand{\RazaoFormatoAcTrinta}{2{,}01\xspace}
\newcommand{\RazaoAcumuladorTakum}{2{,}56\xspace}
\newcommand{\RazaoKernelTrinta}{5{,}15\xspace}
\newcommand{\RazaoAcAmpliadoTrinta}{3{,}15\xspace}
\newcommand{\ParTrintaFPAmbos}{8\xspace}
\newcommand{\ParTrintaFPTakumVence}{7\xspace}
\newcommand{\ParTrintaFPOutroVence}{1\xspace}
\newcommand{\ParTrintaFPEmpates}{0\xspace}
\newcommand{\ParTrintaFPTakumNF}{0\xspace}
\newcommand{\ParTrintaFPOutroNF}{0\xspace}
\newcommand{\ParTrintaAcFPAmbos}{8\xspace}
\newcommand{\ParTrintaAcFPTakumVence}{7\xspace}
\newcommand{\ParTrintaAcFPOutroVence}{1\xspace}
\newcommand{\ParTrintaAcFPEmpates}{0\xspace}
\newcommand{\ParTrintaAcFPTakumNF}{0\xspace}
\newcommand{\ParTrintaAcFPOutroNF}{0\xspace}
\newcommand{\ParDezFPAmbos}{5\xspace}
\newcommand{\ParDezFPTakumVence}{3\xspace}
\newcommand{\ParDezFPOutroVence}{2\xspace}
\newcommand{\ParDezFPEmpates}{0\xspace}
\newcommand{\ParDezFPTakumNF}{0\xspace}
\newcommand{\ParDezFPOutroNF}{3\xspace}
\newcommand{\ParOitoEquatromAmbos}{4\xspace}
\newcommand{\ParOitoEquatromTakumVence}{0\xspace}
\newcommand{\ParOitoEquatromOutroVence}{4\xspace}
\newcommand{\ParOitoEquatromEmpates}{0\xspace}
\newcommand{\ParOitoEquatromTakumNF}{0\xspace}
\newcommand{\ParOitoEquatromOutroNF}{4\xspace}
\newcommand{\ParOitoCincodoisAmbos}{5\xspace}
\newcommand{\ParOitoCincodoisTakumVence}{3\xspace}
\newcommand{\ParOitoCincodoisOutroVence}{2\xspace}
\newcommand{\ParOitoCincodoisEmpates}{0\xspace}
\newcommand{\ParOitoCincodoisTakumNF}{0\xspace}
\newcommand{\ParOitoCincodoisOutroNF}{3\xspace}
\newcommand{\NFTakumOito}{0\xspace}
\newcommand{\NFTakumDez}{0\xspace}
\newcommand{\NFEquatrom}{4\xspace}
\newcommand{\NFCincodois}{3\xspace}
\newcommand{\NFDez}{3\xspace}
\newcommand{\ElemNnzFP}{5{,}33\xspace}
\newcommand{\ElemNnzTakum}{4{,}50\xspace}
\newcommand{\InsnNnzFPCento}{6{,}14\xspace}
\newcommand{\CsrNnzFPCento}{1{,}33\xspace}
\newcommand{\InsnNnzFPDuzentos}{4{,}59\xspace}
\newcommand{\CsrNnzFPDuzentos}{1{,}07\xspace}
\newcommand{\InsnNnzFPQuinhentos}{4{,}34\xspace}
\newcommand{\CsrNnzFPQuinhentos}{1{,}03\xspace}
\newcommand{\InsnNnzTrintaCento}{5{,}98\xspace}
\newcommand{\CsrNnzTrintaCento}{1{,}33\xspace}
\newcommand{\RelInsnTrintaCento}{$-$2{,}7\%\xspace}
\newcommand{\RelCsrTrintaCento}{+0{,}0\%\xspace}
\newcommand{\EconomiaCento}{2{,}7\%\xspace}
\newcommand{\ZvtNnzTrintaCento}{0{,}50\xspace}
\newcommand{\InsnNnzTrintaDuzentos}{4{,}55\xspace}
\newcommand{\CsrNnzTrintaDuzentos}{1{,}07\xspace}
\newcommand{\RelInsnTrintaDuzentos}{$-$0{,}7\%\xspace}
\newcommand{\RelCsrTrintaDuzentos}{+0{,}0\%\xspace}
\newcommand{\EconomiaDuzentos}{0{,}7\%\xspace}
\newcommand{\ZvtNnzTrintaDuzentos}{0{,}37\xspace}
\newcommand{\InsnNnzTrintaQuinhentos}{4{,}33\xspace}
\newcommand{\CsrNnzTrintaQuinhentos}{1{,}03\xspace}
\newcommand{\RelInsnTrintaQuinhentos}{$-$0{,}3\%\xspace}
\newcommand{\RelCsrTrintaQuinhentos}{+0{,}0\%\xspace}
\newcommand{\EconomiaQuinhentos}{0{,}3\%\xspace}
\newcommand{\ZvtNnzTrintaQuinhentos}{0{,}35\xspace}
\newcommand{\InsnNnzDezCento}{4{,}95\xspace}
\newcommand{\CsrNnzDezCento}{1{,}47\xspace}
\newcommand{\RelInsnDezCento}{$-$19{,}3\%\xspace}
\newcommand{\RelCsrDezCento}{+10{,}7\%\xspace}
\newcommand{\ZvtNnzDezCento}{0{,}37\xspace}
\newcommand{\InsnNnzDezDuzentos}{4{,}69\xspace}
\newcommand{\CsrNnzDezDuzentos}{1{,}39\xspace}
\newcommand{\RelInsnDezDuzentos}{+2{,}2\%\xspace}
\newcommand{\RelCsrDezDuzentos}{+29{,}9\%\xspace}
\newcommand{\ZvtNnzDezDuzentos}{0{,}35\xspace}
\newcommand{\InsnNnzDezQuinhentos}{4{,}60\xspace}
\newcommand{\CsrNnzDezQuinhentos}{1{,}36\xspace}
\newcommand{\RelInsnDezQuinhentos}{+6{,}0\%\xspace}
\newcommand{\RelCsrDezQuinhentos}{+32{,}5\%\xspace}
\newcommand{\ZvtNnzDezQuinhentos}{0{,}34\xspace}
\newcommand{\InsnNnzOitoCento}{4{,}51\xspace}
\newcommand{\CsrNnzOitoCento}{1{,}39\xspace}
\newcommand{\RelInsnOitoCento}{$-$26{,}6\%\xspace}
\newcommand{\RelCsrOitoCento}{+4{,}6\%\xspace}
\newcommand{\ZvtNnzOitoCento}{0{,}35\xspace}
\newcommand{\InsnNnzOitoDuzentos}{4{,}43\xspace}
\newcommand{\CsrNnzOitoDuzentos}{1{,}36\xspace}
\newcommand{\RelInsnOitoDuzentos}{$-$3{,}4\%\xspace}
\newcommand{\RelCsrOitoDuzentos}{+27{,}4\%\xspace}
\newcommand{\ZvtNnzOitoDuzentos}{0{,}34\xspace}
\newcommand{\InsnNnzOitoQuinhentos}{4{,}39\xspace}
\newcommand{\CsrNnzOitoQuinhentos}{1{,}35\xspace}
\newcommand{\RelInsnOitoQuinhentos}{+1{,}1\%\xspace}
\newcommand{\RelCsrOitoQuinhentos}{+31{,}2\%\xspace}
\newcommand{\ZvtNnzOitoQuinhentos}{0{,}34\xspace}
\newcommand{\InsnMenorTrinta}{13\xspace}
\newcommand{\InsnIgualTrinta}{11\xspace}
\newcommand{\InsnMaiorTrinta}{0\xspace}
\newcommand{\InsnMenorDez}{11\xspace}
\newcommand{\InsnIgualDez}{0\xspace}
\newcommand{\InsnMaiorDez}{13\xspace}
\newcommand{\InsnMenorOito}{12\xspace}
\newcommand{\InsnIgualOito}{0\xspace}
\newcommand{\InsnMaiorOito}{12\xspace}
\newcommand{\InsnPares}{24\xspace}
\newcommand{\ModRFP}{13\xspace}
\newcommand{\ModCFP}{12\xspace}
\newcommand{\ModKFP}{21\xspace}
\newcommand{\ModROito}{14\xspace}
\newcommand{\ModCOito}{12\xspace}
\newcommand{\ModKOito}{21\xspace}
\newcommand{\ModRDez}{14\xspace}
\newcommand{\ModCDez}{13\xspace}
\newcommand{\ModKDez}{21\xspace}
\newcommand{\ModRTrinta}{14\xspace}
\newcommand{\ModCTrinta}{11\xspace}
\newcommand{\ModKTrinta}{21\xspace}
\newcommand{\ModExatas}{96\xspace}
\newcommand{\ModExec}{96\xspace}
\newcommand{\ModSinteticosOk}{12\xspace}
\newcommand{\ModSinteticosTotal}{12\xspace}
\newcommand{\ModResiduoMax}{0\xspace}


\title[Takum na RVV]{Aritmética Takum na Extensão Vetorial RISC-V}
\subtitle{Projeto Integrador AOC/UFMG}
\author{André Luís Alves Rezende \and Gabriel Ramos Fontes \and Isadora Maria Carvalho Lopes
\and João Pedro Santos de Brito}
\institute{Departamento de Engenharia Eletrônica -- UFMG}
\date{2026}

\begin{document}

\begin{frame}
	\titlepage
\end{frame}

\begin{frame}{Motivação: proliferação de formatos}
	\begin{itemize}
		\item RVV hoje: FP16 (\texttt{Zvfh}), bfloat16
		      (\texttt{Zvfbfmin}/\texttt{Zvfbfwma}), FP32/64, cripto.
		      \textbf{FP8 não existe} na RVV.
		\item Cada formato novo adiciona conversões cruzadas, cantos de exceção
		      e área de decodificador.
		\item A \textbf{aritmética Takum} propõe precisão afunilada com faixa
		      dinâmica limitada e \textbf{um decodificador comum de 12 bits}
		      para todas as larguras.
		\item Replicamos, na RVV, a metodologia do estudo de Hunhold para o
		      AVX10.2 (MOCAST 2025).
	\end{itemize}
\end{frame}

\begin{frame}{Takum em 30 segundos}
	\begin{itemize}
		\item Campos $S, D, R, C, F$; a característica cabe em 12 bits.
		\item T8/T16/T32: códigos curtos usam zeros à direita (janela de 12 bits).
		\item Negação = complemento de dois; NaR único; saturação; sem infinito.
		\item \textbf{Arredondamento RNE no código}: piso no prefixo de $n+1$
		      bits, arredonda pelo bit de meio-ponto e pelo \emph{sticky}.
		      Pode diferir do vizinho mais próximo em valor.
		\item Nesta versão: variante \textbf{linear} (Def.~8).
		      A logarítmica fica para depois ($\Delta\pm$ tem custo próprio).
	\end{itemize}
\end{frame}

\begin{frame}{Zvtakum: cinco instruções, em custom-0}
	\begin{table}
		\centering
		\small
		\begin{tabular}{lll}
			\toprule
			Mnemônico & Operação & Observação \\
			\midrule
			\texttt{vtadd.vv} & $vd_i = vs2_i + vs1_i$ & saturante \\
			\texttt{vtsub.vv} & $vd_i = vs2_i - vs1_i$ & saturante \\
			\texttt{vtmul.vv} & $vd_i = vs2_i \times vs1_i$ & saturante \\
			\texttt{vtdot.vv} & $vd_0 \leftarrow vd_0 + \sum_i vs2_i \times vs1_i$
			                  & acúmulo 2n (T8$\to$T16, \dots) \\
			\texttt{vncvt.t.t} & $T_{2n} \to T_n$ & estreitamento RNE \\
			\bottomrule
		\end{tabular}
	\end{table}
	\begin{itemize}
		\item \texttt{opcode = 0x0B} (custom-0), \texttt{funct3 = 0},
		      \texttt{funct6} = operação, \texttt{vm} no bit 25.
		\item Por que não OP-V: o inventário do riscv-opcodes mostra que
		      \texttt{funct6=0b111001} não é o único livre; custom-0 é o espaço
		      reservado para extensões não padronizadas.
		\item SEW seleciona T8/T16/T32; máscara, cauda e \texttt{vl} da RVV;
		      sem CSRs novos.
	\end{itemize}
\end{frame}

\begin{frame}{Arquitetura: núcleo RISC-V modelado}
	\begin{itemize}
		\item \textbf{Escalar:} RV64IM + Zicsr, com \texttt{write} e \texttt{exit}
		      (ABI Linux), então o kernel roda de verdade no simulador.
		\item \textbf{Vetorial (subconjunto RVV 1.0):} \texttt{vsetvl*},
		      \texttt{vle}/\texttt{vse} (8/16/32), \texttt{vluxei32},
		      \texttt{vmv}, \texttt{vfadd}, \texttt{vfmul},
		      \texttt{vfredusum}; VLEN configurável; LMUL 1--8.
		\item \textbf{Zvtakum:} unidade Takum com decodificador de 12 bits,
		      aritmética exata em 128 bits, acumulador de $2n$ no
		      \texttt{vtdot} e estreitamento no \texttt{vncvt}.
		\item \textbf{Contadores:} instruções por classe, escritas de CSR
		      (inclui \texttt{vl}/\texttt{vtype} do \texttt{vsetvl}),
		      elementos ativos e bytes.
		\item Fora do escopo: modo privilegiado, MMU, cache, FP8 e Zvfh.
	\end{itemize}
\end{frame}

\begin{frame}{Unidade Takum: aritmética própria, sem \texttt{double}}
	\begin{itemize}
		\item Decodificação: janela de 12 bits $\to$ valor exato
		      $(-1)^{S} \cdot \mathrm{sig} \cdot 2^{e-63}$.
		\item Soma/subtração: alinhamento do menor operando com bit
		      \emph{sticky}, normalização e codificação.
		\item Produto: significandos inteiros (até $28\times28$ bits em T32)
		      e soma de expoentes, com 128 bits no acumulador T64.
		\item Codificação RNE no código, com saturação em maxpos/minpos.
		\item \textbf{Sem \texttt{double}} de propósito: as 128 divergências
		      T32 registradas contra a libtakum vêm de arredondamento duplo no
		      intermediário \texttt{double}.
	\end{itemize}
\end{frame}

\begin{frame}{Simulador e verificação}
	\begin{itemize}
		\item \texttt{zt-sim} (C++17, três módulos): \texttt{takum.cpp},
		      \texttt{machine.cpp} e \texttt{main.cpp}. Carrega ELF64 e atende
		      \texttt{write}/\texttt{exit}.
		\item \textbf{Núcleo Takum:} 21 verificações contra vetores dourados
		      gerados pelo modelo Python (T8 completo, narrow T16$\to$T8
		      completo, amostras de T16/T32/T64).
		\item \textbf{Diferencial com QEMU 10.2.2:} escalar e RVV padrão,
		      byte a byte, em VLEN 128/256/512, incluindo máscara com elemento
		      0 inativo, \texttt{ta}/\texttt{ma}, \texttt{vl=0} e
		      \texttt{vstart}.
		\item \textbf{Zvtakum contra o modelo Python:} 1 programa, 3 VLENs,
		      byte a byte.
		\item \textbf{Kernels SpMV:} \NumExecSpmv\ execuções (8 matrizes
		      $\times$ 4 formatos $\times$ 3 VLENs), todas bit a bit iguais ao
		      modelo.
	\end{itemize}
\end{frame}

\begin{frame}{Metodologia experimental}
	\begin{itemize}
		\item \textbf{Matrizes:} 8 da coleção Harwell--Boeing, de 30 a
		      \MaxLinhas\ linhas e até \MaxNnz\ não nulos.
		\item \textbf{Vetor $x$:} pseudoaleatório (semente 2026), elementos
		      $k/2^{52}$: força arredondamento em FP32 e takum32.
		\item \textbf{Kernels:} mesmo laço escalar sobre o CSR; FP32 com
		      \texttt{vfmul}+\texttt{vfredusum}; Takum com
		      \texttt{vtdot}+\texttt{vncvt}; gather \texttt{vluxei32} nos dois.
		\item \textbf{Erro:} relativo L2 contra referência binary128
		      (113 bits, validada contra \texttt{\_\_float128} do GCC).
		\item \textbf{Controles:} mesma matriz e mesmo $x$ com acumuladores de
		      larguras diferentes, para separar formato de acumulação.
		\item \textbf{Comparações pareadas} por matriz: contamos vitórias,
		      derrotas e resultados não finitos.
	\end{itemize}
\end{frame}

\begin{frame}{Resultado 1: precisão em 32 bits}
	\begin{itemize}
		\item Kernel FP32: erro mediano \MedFPtrinta; Takum32: \MedTakumTrinta.
		\item Takum32 vence em \ParTrintaFPTakumVence\ de \ParTrintaFPAmbos\
		      matrizes (exceção: \texttt{pores\_1}).
		\item \textbf{Decomposição:} o ganho de \RazaoKernelTrinta\ se divide em
		      \RazaoFormatoAcTrinta\ do \textbf{formato}, com acumulador de 32
		      bits nos dois, e \RazaoAcumuladorTakum\ do \textbf{acumulador} de
		      64 bits do \texttt{vtdot}.
		\item Com o mesmo acumulador de 32 bits, o Takum32 vence em
		      \ParTrintaAcFPTakumVence\ de \ParTrintaAcFPAmbos\ matrizes.
		\item Com acumulador de 64 bits nos dois, Takum32 é
		      \RazaoAcAmpliadoTrinta\ vezes melhor.
	\end{itemize}
\end{frame}

\begin{frame}{Resultado 2: 8 e 16 bits}
	\begin{itemize}
		\item Takum8 e Takum16 \textbf{não estouram} em nenhuma matriz.
		      E4M3 estoura em \NFEquatrom, E5M2 em \NFCincodois\ e FP16 em
		      \NFDez\ das 8.
		\item Onde nenhum estoura:
		\begin{itemize}
			\item Takum16 vence FP16 (acumulador 32 bits) em
			      \ParDezFPTakumVence\ de \ParDezFPAmbos;
			\item Takum8 vence E5M2 (acumulador 16 bits) em
			      \ParOitoCincodoisTakumVence\ de \ParOitoCincodoisAmbos;
			\item E4M3 vence Takum8 em \ParOitoEquatromOutroVence\ de
			      \ParOitoEquatromAmbos.
		\end{itemize}
		\item Leitura honesta: o Takum troca precisão por ausência de overflow.
	\end{itemize}
\end{frame}

\begin{frame}{Resultado 3: custo de instruções}
	\begin{itemize}
		\item Modelo analítico: $N = rR + cS + K$ ($R$ = linhas, $S$ = faixas),
		      com parâmetros medidos e \textbf{resíduo zero} nas \ModExec\
		      execuções.
		\item Diferença do Takum32 para o FP32: $N_{\text{Takum32}} -
		      N_{\text{FP32}} = R - S$. Um \texttt{vncvt} a mais por linha,
		      um \texttt{vtdot} a menos por faixa.
		\item O Takum32 nunca executa mais que o FP32:
		      \InsnMenorTrinta\ pares com menos, \InsnIgualTrinta\ iguais,
		      \InsnMaiorTrinta\ com mais. Escritas de CSR: idênticas.
		\item Em VLEN 256: FP32 \InsnNnzFPDuzentos\ instruções por não nulo,
		      Takum32 \InsnNnzTrintaDuzentos\ ($\RelInsnTrintaDuzentos$).
		\item O de 8 e 16 bits paga um \texttt{vsetvli} extra por faixa
		      (índice de 32 bits com LMUL maior): \RelCsrOitoDuzentos\ e
		      \RelCsrDezDuzentos\ mais escritas de CSR em VLEN 256.
	\end{itemize}
\end{frame}

\begin{frame}{Limitações e trabalhos futuros}
	\begin{itemize}
		\item 8 matrizes e um único vetor $x$: ordens de grandeza, não
		      intervalos de confiança.
		\item Estudo de 1401 matrizes (MuFoLAB) não executado.
		\item Variante logarítmica não implementada; FP16 vetorial (Zvfh)
		      apenas reconhecida pelo montador.
		\item Contadores medem instruções, não ciclos, área ou energia.
		\item Futuro: 1401 matrizes, Zvfh, $\Delta\pm$, kernel sem a troca de
		      configuração do índice e comparação com hardware.
	\end{itemize}
\end{frame}

\begin{frame}{Conclusão}
	\begin{itemize}
		\item Extensão Zvtakum especificada e implementada em custom-0, com
		      aritmética Takum própria em inteiros.
		\item Simulador funcional RV64IM+RVV+Zvtakum verificado contra o QEMU
		      (byte a byte) e contra um modelo de racionais exatos.
		\item Takum32: erro menor que o FP32 em \ParTrintaAcFPTakumVence\ de
		      \NumMatrizes\ matrizes com acumulador de mesma largura, e nunca
		      mais instruções.
		\item Em 8 e 16 bits o custo depende do VLEN; o Takum não estoura,
		      mas não é sempre mais preciso.
	\end{itemize}
	\vfill
	\centering \large Obrigado! Perguntas?
\end{frame}

\begin{frame}{Demo ao vivo}
	\begin{itemize}
		\item \texttt{make sim-test} --- núcleo Takum contra vetores dourados.
		\item \texttt{make sim-check} --- diferencial com QEMU e com o modelo
		      (9 execuções, byte a byte).
		\item Execução de um kernel no simulador, com contadores:\\
		      \texttt{sim/build/zt-sim
		      sim/build/spmv/bcsstk01.takum32.v256.elf
		      -{}-vlen 256 -{}-stdout /tmp/y.bin
		      -{}-counters /tmp/c.json}
		\item \texttt{make spmv} --- 96 kernels + FP8/FP16 + modelo de
		      contagem (regenera os CSVs do artigo).
	\end{itemize}
\end{frame}

\end{document}
